% !TEX root = ../main.tex

\chapter{实验设计与结果分析}

\section{实验目标与实验思路}

本文实验主要验证两个问题：监督微调是否能提升基础大语言模型在股票短期涨跌方向判断上的表现；在预测输入中加入新闻和研报等外部文本信息后，模型是否会比仅依赖历史 K 线时表现更好。围绕这两个问题，本文设计了“微调前后对比实验”和“有无外部信息对比实验”。

当前 SerpAPI 模块主要服务于实时网站展示，用于在线查询时补充搜索引擎返回的最新网页结果。由于历史时点的通用搜索结果难以稳定回放和严格复现，本文离线实验暂未将 SerpAPI 返回结果纳入统一评测数据，而是先使用可稳定回放的 news 与 report 摘要完成方法验证。

实验采用离线历史回放式验证。构造每个测试样本时，系统只使用目标日期之前或当天可见的信息，包括历史行情、新闻摘要和研报摘要；随后将模型预测结果与未来真实涨跌方向比较。该方法不能完全等同于真实交易环境，但可以在避免未来信息泄漏的前提下，对模型方向判断能力进行可复现评估。

股票日频收益噪声较高，短期涨跌幅本身也很难精确预测，因此本文不直接回归未来收益率，而是采用涨跌方向分类。7 日 K 线任务只要求模型输出下一交易日“涨”或“跌”；多源信息任务则要求同时输出未来 1、3、7 个交易日后的方向组合。该设置便于与真实标签比较，也更接近投研中先判断方向、再讨论仓位和风险控制的流程。

\section{数据集与离线验证方案}

\subsection{历史样本来源}

实验数据主要来自历史行情、新闻和研报采集结果。行情数据包含股票代码、交易日期、开盘价、最高价、最低价、收盘价、成交量、成交额、振幅、涨跌幅和换手率等字段。新闻和研报数据则进一步配有正文或摘要缓存。

对于监督微调主线，本文构造了 7 日 K 线涨跌方向数据集。该数据集以连续 7 个交易日的归一化 K 线作为输入，以第 8 个交易日的真实涨跌方向作为标签。当前该数据集包含 339287 条训练样本，字段为 \texttt{prompt} 和 \texttt{label}。

为了训练模型生成带分析过程的回答，本文进一步构建了 CoT 风格数据。该过程基于基础 7 日 K 线样本和真实标签生成批处理请求，调用 \texttt{qwen3-max} 为每条样本生成解释性推理，并要求最终答案与真实标签一致；任务完成后，再将批处理结果与原始样本按 \texttt{custom\_id} 对齐，合并为 \texttt{prompt} 和 \texttt{completion} 两列。当前该数据集包含 599871 条样本，用于 LoRA SFT。

对于多源信息验证，本文构造了包含 K 线、新闻摘要和研报摘要的样本。该数据集以目标日期及之前连续 5 个交易日的归一化 K 线、最近 30 日内的新闻摘要、最近 30 日内的研报摘要作为输入，以未来 1、3、7 个交易日相对目标日期收盘价的涨跌方向作为标签。当前该实验共包含 1691 条样本。

\subsection{离线历史回放式验证方法}

离线验证的关键是保证样本输入不包含未来信息。对于行情特征，训练和测试样本只使用目标日期之前或当天的 K 线数据。对于新闻和研报摘要，系统只选择发布日期不晚于目标日期的内容，并额外限制新闻和研报的最大回看窗口，避免模型看到目标日期之后才出现的信息。

在 7 日 K 线任务中，训练样本默认使用 2021 年 1 月 1 日至 2026 年 1 月 1 日之前的标签日，评测样本使用 2026 年 1 月 1 日至 2026 年 3 月 1 日之前的标签日。评测过程会分别调用基础模型和加载 LoRA adapter 后的微调模型，对同一批 eval prompts 进行生成，并从输出中解析最终“涨”或“跌”标签。

在多源信息任务中，评测过程将每条样本组织为包含 K 线、news 摘要和 report 摘要的完整 Prompt，并要求模型输出未来 1、3、7 个交易日的方向组合。为了构造“无外部信息”对照组，系统使用同一数据集中的 K 线和标签，但在 Prompt 中移除新闻与研报摘要，仅保留历史行情信息。

\subsection{训练集、验证集与测试集划分}

监督微调训练集主要来自 2026 年以前的历史行情样本。若未提供独立验证集，系统会从训练数据中按默认比例划分少量验证样本，用于训练过程中监控 eval loss。本文关注的是最终离线测试表现，因此主要报告独立时间窗口上的测试结果。

7 日 K 线评测集的标签日窗口为 2026 年 1 月 1 日至 2026 年 3 月 1 日之前，样本数为 1330。该数据集字段为 \texttt{prompt} 和 \texttt{completion}，其中 \texttt{completion} 中只保留真实最终方向标签。多源信息评测集的目标日期范围为 2026 年 1 月 1 日至 2026 年 4 月 1 日，样本数为 1691，其标签为未来 1、3、7 个交易日的方向组合。

\section{实验设置}

\subsection{对比模型设置}

本文主要对比以下模型或输入设置。

基础模型为 Qwen2.5-7B-Instruct。该模型不经过本文任务数据微调，直接根据 Prompt 生成涨跌方向预测，用于衡量通用指令模型在该任务上的原始能力。

SFT LoRA 微调模型以 Qwen2.5-7B-Instruct 为基座，在 CoT 风格股票方向预测数据上进行 LoRA 监督微调。LoRA 默认设置为 \(r=64\)、\(\alpha=128\)、dropout 为 0.05，训练过程中只更新 adapter 参数。

仅行情输入设置只向模型提供归一化 K 线信息，不提供新闻或研报摘要，用作外部信息增强效果的基线。

多源信息输入设置在 K 线之外加入目标日期之前可见的 news 摘要和 report 摘要，用于验证外部文本信息对未来方向预测的帮助。SerpAPI 所代表的 Web 检索结果当前不进入该离线评测集合，而是作为在线系统中的实时补充模块单独实现。

\subsection{有无外部信息实验设置}

有无外部信息实验基于同一批多周期样本进行。仅行情组的输入只包含目标日期及之前连续 5 个交易日的归一化 K 线；外部信息组在相同 K 线特征之外，额外加入新闻摘要和研报摘要。

当前外部信息组评测使用的是加载 LoRA adapter 后的模型，因此该结果更接近“完整系统设置”的验证，而不是严格控制模型参数完全相同的单因素消融实验。严格的同模型外部信息消融可以作为后续实验补充。

\subsection{不同预测周期实验设置}

多源信息实验同时评估未来 1、3、7 个交易日三个预测周期。标签生成方式为：分别比较目标日期收盘价与未来第 1、第 3、第 7 个交易日收盘价，若未来收盘价大于或等于目标日期收盘价，则记为“涨”，否则记为“跌”。模型输出需要按照固定格式给出三个方向，例如“涨，跌，涨”。

为了更细粒度地分析不同周期的预测效果，本文除了计算三项平均逐项准确率外，还统计每个预测周期上的单独准确率，以及三项完全匹配率。前者反映模型在每个时间尺度上的平均方向判断能力，后者要求三个预测周期全部正确，因此更加严格。

\subsection{评价指标设置}

对于单周期涨跌二分类任务，本文采用准确率作为评价指标：

\[
\mathrm{Acc}=\frac{N_{\mathrm{correct}}}{N_{\mathrm{total}}},
\]

其中 \(N_{\mathrm{correct}}\) 表示预测方向与真实标签一致的样本数，\(N_{\mathrm{total}}\) 表示参与评估的样本总数。

对于 1、3、7 日多周期任务，本文采用平均逐项准确率。若第 \(i\) 个样本包含三个预测位置，则其得分为三个位置中预测正确比例，整体准确率为所有样本得分的平均值。若模型输出无法解析为合法的三项涨跌标签，则记为提取失败；在包含提取失败的统计中，提取失败样本按 0 分计入，在排除提取失败的统计中则从分母中剔除。

\section{微调前后模型性能对比分析}

\subsection{无额外信息条件下的结果分析}

表 \ref{tab:sft-compare} 给出了 7 日 K 线单周期任务中基础模型与 SFT LoRA 模型的对比结果。该实验只使用历史行情特征，不引入新闻或研报等外部文本信息。

\begin{table}[htbp]
  \centering
  \caption{7 日 K 线单周期任务中微调前后模型结果}
  \label{tab:sft-compare}
  \begin{tabular}{lccc}
    \toprule
    模型 & 正确数 & 样本数 & 准确率 \\
    \midrule
    Qwen2.5-7B-Instruct & 658 & 1330 & 49.47\% \\
    Qwen2.5-7B-Instruct + LoRA SFT & 694 & 1330 & 52.18\% \\
    \bottomrule
  \end{tabular}
\end{table}

基础模型在该任务上的准确率为 49.47\%，接近随机二分类水平；经过 LoRA 监督微调后，准确率提升至 52.18\%，提高 2.71 个百分点。这个变化不算大，说明仅依赖历史 K 线预测下一交易日涨跌仍然较难；但在相同评测窗口下，LoRA SFT 确实让模型更适应该任务的输入格式和输出约束。

从实验意义上看，SFT 的主要作用不只是提升数值准确率，还包括增强模型输出的可解析性。评测过程需要从模型回复中提取最终“涨”或“跌”标签，若模型不遵循格式，则会影响自动评估。当前记录中两组模型均未出现提取失败样本，说明 Prompt 约束和 CoT 风格训练数据能够较好地维持输出格式稳定。

\subsection{有额外信息条件下的结果分析}

表 \ref{tab:external-compare} 展示了多周期任务中仅行情输入与加入外部信息后的结果。该任务要求模型同时预测未来 1、3、7 个交易日的方向，因此表中准确率采用平均逐项准确率。

\begin{table}[htbp]
  \centering
  \caption{多周期任务中有无外部信息的结果对比}
  \label{tab:external-compare}
  \begin{tabular}{lcccc}
    \toprule
    设置 & 样本数 & 提取失败数 & 含失败准确率 & 排除失败准确率 \\
    \midrule
    仅 K 线输入 & 1691 & 18 & 49.85\% & 50.39\% \\
    K 线 + news/report 摘要 & 1691 & 97 & 52.91\% & 56.13\% \\
    \bottomrule
  \end{tabular}
\end{table}

在包含提取失败的统计下，加入 news 和 report 摘要后的平均逐项准确率为 52.91\%，仅 K 线输入为 49.85\%，提升 3.06 个百分点；排除提取失败后，外部信息组准确率达到 56.13\%，仅 K 线组为 50.39\%，提升 5.74 个百分点。这个差异说明，新闻和研报在部分样本中提供了行情数据之外的事件解释和基本面线索。这一现象与新闻文本、金融文本融合和大语言模型股票预测相关研究的结论基本一致\cite{lopezlira2023chatgptstock,elahi2024financialnewsllm}。

外部信息组的提取失败数也更多。输入变长、信息变复杂后，模型更容易出现输出格式不完全符合要求的情况。多源信息增强虽然提升了有效样本上的判断能力，但也提高了 Prompt 设计、输出约束和结果解析的要求。

\subsection{不同预测周期结果分析}

表 \ref{tab:horizon-compare} 给出了有效样本上的分周期准确率。仅行情输入组在 1 日和 3 日预测上的准确率分别为 45.49\% 和 41.18\%，加入外部信息后分别提升到 52.38\% 和 52.26\%，对应提升 6.89 和 11.08 个百分点。news 和 report 摘要对较短周期方向判断帮助更明显，可能是因为短期走势更容易受到事件、资金情绪和市场叙事影响。这一现象也与社交媒体情绪和新闻事件对短期市场反应具有解释力的相关研究相符\cite{bollen2011twitter,lopezlira2023chatgptstock}。

\begin{table}[htbp]
  \centering
  \caption{不同预测周期上的准确率对比（排除提取失败）}
  \label{tab:horizon-compare}
  \begin{tabular}{lccc}
    \toprule
    设置 & 1 个交易日 & 3 个交易日 & 7 个交易日 \\
    \midrule
    仅 K 线输入 & 45.49\% & 41.18\% & 64.49\% \\
    K 线 + news/report 摘要 & 52.38\% & 52.26\% & 63.74\% \\
    \bottomrule
  \end{tabular}
\end{table}

7 日预测上，仅 K 线输入组为 64.49\%，外部信息组为 63.74\%，下降 0.75 个百分点。较长周期的方向标签受更多后续市场因素影响，当前日期附近的新闻和研报摘要不一定能稳定解释未来一周走势；该结果也可能与样本分布、标签不平衡或模型输出偏好有关。从当前数据看，外部信息的提升主要集中在 1 日和 3 日预测上。

三项完全匹配率上，外部信息组为 21.82\%，仅行情输入组为 11.12\%，提高 10.70 个百分点。完全匹配要求三个预测周期全部正确，本身比逐项准确率更严格。外部信息组在该指标上的提升说明，多源文本不仅提高了单个位置的平均准确率，也改善了部分样本上的多周期方向组合判断。

\section{强化学习优化实验设计}

本文当前实验结果主要来自基础模型、LoRA SFT 模型以及有无外部信息输入的离线评测。强化学习部分已经完成训练入口和奖励设计，但尚未形成完整的 SFT+RL 结果表。因此，本节不报告强化学习后的准确率提升，只说明后续实验设计，避免把尚未完成的工作写成已验证结论。

在模型对比方面，后续实验应在相同测试集上比较 SFT 模型与 SFT+RL 模型。SFT 模型作为初始策略，已经具备基本的任务格式遵循能力；RL 模型则在此基础上，基于纯 7 日 K 线构造的多周期方向样本，通过未来 1、3、7 个交易日方向逐项平均得分、输出格式稳定性和理由一致性等奖励信号进一步优化。评价指标仍可沿用本文的单周期准确率、多周期平均逐项准确率、三项完全匹配率和提取失败率。

在与外部信息实验的衔接方面，后续可以先在纯 K 线多周期任务上验证 RL 是否能够稳定提升基础方向判断能力，再进一步考察这种优化能否迁移到包含 news/report 摘要的多源输入场景中。如果 RL 只提升纯 K 线条件下的格式遵循，而无法改善多源条件下的有效样本准确率，则说明当前奖励设计更偏向表层格式；如果 1 日和 3 日预测在多源任务中也进一步改善，再讨论奖励信号对短期信息利用的帮助会更有依据。

在训练稳定性方面，后续需要重点记录奖励曲线、KL 距离、输出长度、提取失败率和不同周期准确率的变化。股票预测任务噪声较强，如果奖励设计过于单一，模型可能过度拟合某一类标签或输出模板。因此，RL 实验应当设置格式惩罚、长度约束和相对 SFT 模型的稳定性约束，使优化目标集中在预测质量和证据一致性上，而不是牺牲可读性与泛化能力来追求局部得分。

\section{实验结果讨论}

\subsection{实验结果解释}

7 日 K 线单周期任务中，SFT 后准确率从 49.47\% 提升到 52.18\%，提高 2.71 个百分点。这个提升幅度并不大，但它发生在同一评测集和同一解析规则下，说明微调主要改善了模型对任务格式和行情输入的适配。考虑到股票日频方向本身接近高噪声二分类，这一结果不能被解读为稳定预测能力。

多周期实验中，news 和 report 摘要带来的增益主要集中在 1 日和 3 日预测上，7 日预测差异并不明显。这一结果符合直觉：新闻和研报摘要通常更接近近期事件或短期预期，对一周之后的价格方向未必仍有稳定解释力。SerpAPI 已经接入在线系统，但由于搜索结果难以做历史回放，本节暂时只把可复现的 news/report 摘要纳入离线比较。

\subsection{方法局限性分析}

实验也有明显限制。当前外部信息组与仅行情组还不是完全严格的单因素消融，部分设置同时包含模型微调和输入内容变化，因此不能把全部提升都归因于外部文本。更严格的做法是在同一模型和同一解码参数下，分别比较仅 K 线、K 线加 news、K 线加 report 和完整输入。

另一个问题是输出解析。多源输入变长后，模型更容易偏离固定格式，提取失败样本随之增加。后续如果继续做量化评测，应优先把输出改成 JSON schema、函数调用式输出或分类头，而不是只依赖自然语言 Prompt。当前样本覆盖的股票池和时间区间也有限，结果还不能外推到所有市场阶段。

\section{本章小结}

本章的实验主要回答两个问题：微调是否有帮助，外部文本是否有帮助。结果显示，LoRA SFT 在 7 日 K 线单周期任务上提升 2.71 个百分点；加入 news/report 摘要后，1 日和 3 日方向预测分别提升 6.89 和 11.08 个百分点，但 7 日预测没有提升。与此同时，实验也暴露出两个后续必须处理的问题：多源输入更容易带来格式解析失败，SerpAPI 这类实时检索结果还缺少可回放评测数据。因此，本章结果更适合作为方法可行性验证，而不是最终性能结论。
